\documentclass[11pt]{article}
\usepackage[margin=1in]{geometry}
\usepackage{amsmath,amssymb,bm}
\usepackage[T1]{fontenc}
\usepackage[utf8]{inputenc}
\usepackage{lmodern}
\usepackage{booktabs}
\usepackage{longtable}
\usepackage{array}
\usepackage{enumitem}
\usepackage{hyperref}

\title{Augmented ESKF Mathematical Formulation}
\author{sensor-fusion documentation}
\date{\today}

\newcommand{\R}{\mathbb{R}}
\newcommand{\T}{^{T}}
\newcommand{\crossmat}[1]{\left[#1\right]_\times}
\newcommand{\dq}{\delta q}
\newcommand{\dt}{\Delta t}

\begin{document}
\maketitle
\tableofcontents
\newpage

\section{Scope}

This document describes the runtime augmented ESKF implemented by
\texttt{ekf/src/rust\_eskf.rs}, \texttt{ekf/src/generated\_eskf.rs}, and the
included generated model files. The filter estimates NED attitude, velocity,
position, IMU biases, and a residual seeded-frame-to-vehicle mount quaternion.
It is the primary runtime inertial/GNSS filter behind \texttt{SensorFusion}.

\section{Frames and Mount Sources}

\begin{longtable}{>{\raggedright\arraybackslash}p{0.12\linewidth} >{\raggedright\arraybackslash}p{0.78\linewidth}}
\toprule
\textbf{Symbol} & \textbf{Meaning} \\
\midrule
\(n\) & Local NED frame. ESKF position and velocity states are in this frame. \\
\(b\) & Raw IMU body frame. \\
\(v\) & Physical vehicle frame, FRD: forward, right, down. \\
\(s\) & Seeded vehicle frame after raw IMU samples are pre-rotated by a coarse mount seed. \\
\(c\) & Corrected vehicle frame after applying the residual mount quaternion \(q_{cs}\). \\
\bottomrule
\end{longtable}

Rotations follow
\begin{equation}
  x_a=C_{ab}x_b,\qquad C(q_1\otimes q_2)=C(q_1)C(q_2).
\end{equation}
The align filter supplies a physical vehicle-to-body mount \(q_{vb}^{0}\). The
fusion facade pre-rotates each raw IMU delta into the seeded frame:
\begin{equation}
  \Delta\theta_s=C_{sv}^{0}C_{vb}^{0}\Delta\theta_b
  =C_{bv}^{0\T}\Delta\theta_b,
  \qquad
  \Delta v_s=C_{bv}^{0\T}\Delta v_b .
\end{equation}
If \texttt{EskfMountSource::LatchedAlign} is selected, \(q_{vb}^{0}\) is frozen
at ESKF initialization. If \texttt{FollowAlign} is selected, the pre-rotation
source follows the current align mount and the residual mount is reset to
identity. A fixed external mount is also supported.

\section{Nominal and Error States}

The nominal state has 20 scalar components:
\begin{equation}
  x=
  \begin{bmatrix}
  q_{ns} & v_n\T & p_n\T & b_g\T & b_a\T & q_{cs}
  \end{bmatrix}\T ,
\end{equation}
where \(q_{ns}\) rotates seeded-frame vectors into NED, \(v_n\) is NED
velocity, \(p_n\) is local NED position, \(b_g\) and \(b_a\) are additive gyro
and accelerometer biases in the seeded frame, and \(q_{cs}\) rotates
seeded-frame velocity into the corrected vehicle frame.

The 18-dimensional error state is
\begin{equation}
  \delta x =
  \begin{bmatrix}
  \delta\theta_s &
  \delta v_n &
  \delta p_n &
  \delta b_g &
  \delta b_a &
  \delta\psi_{cs}
  \end{bmatrix}\T .
\end{equation}
The Rust index order is:
\begin{center}
\begin{tabular}{ll}
\toprule
Indices & Error component \\
\midrule
0..2 & attitude small angle \(\delta\theta_s\) \\
3..5 & velocity error \(\delta v_n\) \\
6..8 & position error \(\delta p_n\) \\
9..11 & gyro-bias error \(\delta b_g\) \\
12..14 & accelerometer-bias error \(\delta b_a\) \\
15..17 & residual mount error \(\delta\psi_{cs}\) \\
\bottomrule
\end{tabular}
\end{center}

\section{Nominal Propagation}

The ESKF consumes one-sample IMU increments already expressed in the seeded
frame:
\begin{equation}
  u_k = \{\Delta\theta_s,\Delta v_s,\dt\}.
\end{equation}
Bias-corrected increments are
\begin{equation}
  \Delta\theta_s^c=\Delta\theta_s-b_g\dt,\qquad
  \Delta v_s^c=\Delta v_s-b_a\dt .
\end{equation}
The generated nominal equations implement
\begin{align}
  q_{ns}^{+} &= \operatorname{normalize}\left(q_{ns}^{-}\otimes
  \dq(\Delta\theta_s^c)\right),\\
  v_n^{+} &= v_n^{-}+C_{ns}(q_{ns}^{+})\Delta v_s^c+g_n\dt,\\
  p_n^{+} &= p_n^{-}+v_n^{+}\dt,\\
  b_g^{+}&=b_g^{-},\qquad b_a^{+}=b_a^{-},\qquad q_{cs}^{+}=q_{cs}^{-},
\end{align}
with \(g_n=[0,0,9.80665]\T\) in NED.

\section{Error-State Propagation}

The covariance prediction uses generated discrete matrices:
\begin{equation}
  \delta x_{k+1}=F_k\delta x_k+G_kw_k,
  \qquad
  P_{k+1}^{-}=F_kP_k^{+}F_k\T+G_kQ_kG_k\T .
\end{equation}
The 15-noise vector is
\begin{equation}
  w=
  \begin{bmatrix}
  n_g & n_a & n_{bg} & n_{ba} & n_m
  \end{bmatrix}\T .
\end{equation}
The diagonal covariance entries assembled at runtime are
\begin{equation}
  Q_k=\operatorname{diag}\left(
  \sigma_g^2\dt^2 I_3,\,
  \sigma_a^2\dt^2 I_3,\,
  \sigma_{bg}^2\dt I_3,\,
  \sigma_{ba}^2\dt I_3,\,
  \sigma_m^2\dt I_3
  \right).
\end{equation}
If residual mount states are frozen, the mount random-walk block is set to zero
and the mount covariance rows/columns are cleared.

\section{Measurement Updates}

Scalar observations use
\begin{align}
  r &= z-h(\hat x),\\
  S &= HPH\T+R,\\
  K &= PH\T S^{-1},\\
  \delta x &= Kr .
\end{align}
Covariance is updated with a scalar Joseph-form expression. The lateral and
vertical NHC rows may be fused as one sequential two-row batch so the second
row sees the covariance and accumulated correction from the first row.

\subsection{GNSS Position and Velocity}

GNSS position and velocity are local NED:
\begin{align}
  h_p(x)&=p_n,& H_p&=\begin{bmatrix}0&0&I_3&0&0&0\end{bmatrix},\\
  h_v(x)&=v_n,& H_v&=\begin{bmatrix}0&I_3&0&0&0&0\end{bmatrix}.
\end{align}
The generated GNSS rows have no direct residual-mount Jacobian. The runtime has
optional mount update scales for GNSS rows; with the current generated \(H\)
these only affect mount updates induced through covariance cross-correlation.
The default public \texttt{fuse\_gps} path passes zero scales.

\subsection{Zero Velocity}

The zero-velocity pseudo-measurement is a GNSS-velocity-shaped update:
\begin{equation}
  z_v=0,\qquad h_v(x)=v_n .
\end{equation}
It does not directly observe residual mount because it is expressed in NED, not
in corrected vehicle coordinates.

\subsection{Vehicle Speed and NHC}

Predicted seeded-frame velocity is
\begin{equation}
  v_s=C_{ns}\T v_n .
\end{equation}
The corrected vehicle-frame velocity is
\begin{equation}
  v_c=C_{cs}v_s .
\end{equation}
The optional wheel/CAN speed measurement observes
\begin{equation}
  h_x(x)=e_x\T v_c .
\end{equation}
The non-holonomic constraints use zero lateral and vertical vehicle velocity:
\begin{equation}
  h_y(x)=e_y\T v_c\approx 0,\qquad
  h_z(x)=e_z\T v_c\approx 0 .
\end{equation}
For a left residual-mount perturbation,
\begin{equation}
  q_{cs}^{true}=\dq(\delta\psi_{cs})\otimes q_{cs},
\end{equation}
the vehicle velocity perturbation is
\begin{equation}
  v_c^{true}\simeq (I+\crossmat{\delta\psi_{cs}})v_c
  = v_c-\crossmat{v_c}\delta\psi_{cs}.
\end{equation}
Therefore the residual-mount Jacobian block for component \(i\) is
\begin{equation}
  H_{\psi,i}=-e_i\T\crossmat{v_c}.
\end{equation}
Runtime controls can ramp, gate, attenuate, or freeze the residual-mount part
of the Kalman gain and injected correction for these rows.

\subsection{Stationary Gravity}

Stationary gravity updates use body/seed-frame accelerometer samples and update
only \(x\) and \(y\) accelerometer channels:
\begin{equation}
  f_s \approx -C_{ns}\T g_n + b_a+\eta .
\end{equation}
The scalar rows are
\begin{equation}
  h_x=e_x\T(-C_{ns}\T g_n+b_a),\qquad
  h_y=e_y\T(-C_{ns}\T g_n+b_a).
\end{equation}
Before each stationary gravity scalar update, the runtime floors roll and pitch
covariance to avoid over-confident attitude locking. The residual-mount
Jacobian block is zero because the accelerometer sample has already been put in
the seeded frame.

\section{Injection and Covariance Reset}

After each update the error state is injected:
\begin{align}
  q_{ns}^{+} &= q_{ns}^{-}\otimes\dq(\delta\theta_s),\\
  v_n^{+} &= v_n^{-}+\delta v_n,\qquad
  p_n^{+}=p_n^{-}+\delta p_n,\\
  b_g^{+} &= b_g^{-}+\delta b_g,\qquad
  b_a^{+}=b_a^{-}+\delta b_a,\\
  q_{cs}^{+} &= \dq(\delta\psi_{cs})\otimes q_{cs}^{-}.
\end{align}
Both quaternions are normalized. The attitude and residual-mount covariance
blocks are reset with
\begin{equation}
  G_r(\delta\theta)\simeq I-\frac{1}{2}\crossmat{\delta\theta}.
\end{equation}
The full reset matrix is block diagonal with \(G_r\) on the attitude and mount
blocks and identity elsewhere:
\begin{equation}
  P^{+}\leftarrow G_{\mathrm{reset}}P^{+}G_{\mathrm{reset}}\T .
\end{equation}

\section{Initialization and Mount-Source Modes}

At GNSS initialization the filter sets \(q_{ns}\) from GNSS heading or horizontal
velocity course, copies GNSS NED position and velocity into \(p_n,v_n\), sets
biases to zero, and initializes \(q_{cs}=1\). Covariance is diagonal with
configuration-controlled bias and residual-mount variances.

The fusion facade controls the coarse mount used for IMU pre-rotation:
\begin{itemize}[leftmargin=1.5em]
\item \textbf{Latched/internal}: align produces \(q_{vb}^{0}\), which is frozen
when the ESKF initializes. Residual mount \(q_{cs}\) may then estimate remaining
seed error.
\item \textbf{FollowAlign}: the pre-rotation source tracks align and the
residual \(q_{cs}\) is repeatedly reset to identity. This is useful when the
ESKF should not own mount estimation.
\item \textbf{External/reference}: a fixed supplied mount or reference mount is
used as the pre-rotation source.
\end{itemize}

\section{Known Modeling Boundaries}

\begin{itemize}[leftmargin=1.5em]
\item Residual mount does not affect inertial propagation; it is observable only
through vehicle-frame velocity measurements and through cross-covariances.
\item GNSS position, GNSS velocity, zero velocity, and stationary gravity do not
directly observe \(q_{cs}\).
\item Lateral and vertical NHC observations can couple attitude, velocity,
biases, and residual mount; aggressive mount updates are therefore gated and
ramped in the runtime facade.
\item The physical full mount used for plotting must account for the frozen
coarse seed and the direction of \(q_{cs}\). Under the convention above, a
vehicle-to-body comparison uses the seed composed with the inverse residual,
not the residual alone.
\end{itemize}

\end{document}
